% Options for packages loaded elsewhere
\PassOptionsToPackage{unicode}{hyperref}
\PassOptionsToPackage{hyphens}{url}
\PassOptionsToPackage{dvipsnames,svgnames,x11names}{xcolor}
%
\documentclass[
  10pt,
]{article}
\usepackage{amsmath,amssymb}
\usepackage{lmodern}
\usepackage{iftex}
\ifPDFTeX
  \usepackage[T1]{fontenc}
  \usepackage[utf8]{inputenc}
  \usepackage{textcomp} % provide euro and other symbols
\else % if luatex or xetex
  \usepackage{unicode-math}
  \defaultfontfeatures{Scale=MatchLowercase}
  \defaultfontfeatures[\rmfamily]{Ligatures=TeX,Scale=1}
  \setmainfont[]{DejaVu Serif}
  \setmonofont[]{DejaVu Sans Mono}
\fi
% Use upquote if available, for straight quotes in verbatim environments
\IfFileExists{upquote.sty}{\usepackage{upquote}}{}
\IfFileExists{microtype.sty}{% use microtype if available
  \usepackage[]{microtype}
  \UseMicrotypeSet[protrusion]{basicmath} % disable protrusion for tt fonts
}{}
\makeatletter
\@ifundefined{KOMAClassName}{% if non-KOMA class
  \IfFileExists{parskip.sty}{%
    \usepackage{parskip}
  }{% else
    \setlength{\parindent}{0pt}
    \setlength{\parskip}{6pt plus 2pt minus 1pt}}
}{% if KOMA class
  \KOMAoptions{parskip=half}}
\makeatother
\usepackage{xcolor}
\IfFileExists{xurl.sty}{\usepackage{xurl}}{} % add URL line breaks if available
\IfFileExists{bookmark.sty}{\usepackage{bookmark}}{\usepackage{hyperref}}
\hypersetup{
  colorlinks=true,
  linkcolor={blue},
  filecolor={Maroon},
  citecolor={Blue},
  urlcolor={blue},
  pdfcreator={LaTeX via pandoc}}
\urlstyle{same} % disable monospaced font for URLs
\usepackage[margin=1in]{geometry}
\setlength{\emergencystretch}{3em} % prevent overfull lines
\providecommand{\tightlist}{%
  \setlength{\itemsep}{0pt}\setlength{\parskip}{0pt}}
\setcounter{secnumdepth}{-\maxdimen} % remove section numbering
\ifLuaTeX
  \usepackage{selnolig}  % disable illegal ligatures
\fi

\author{}
\date{}

\begin{document}

\hypertarget{falsification-first-decipherment-a-decontaminated-inference-framework-for-testing-undeciphered-script-claims-with-linear-a-as-a-worked-null}{%
\section{Falsification-First Decipherment: A Decontaminated Inference
Framework for Testing Undeciphered-Script Claims, with Linear A as a
Worked
Null}\label{falsification-first-decipherment-a-decontaminated-inference-framework-for-testing-undeciphered-script-claims-with-linear-a-as-a-worked-null}}

Kyriakos Papadopoulos

\hypertarget{abstract}{%
\subsection{Abstract}\label{abstract}}

Decipherment research has a false-positive problem and no shared
instrument for measuring it: a small corpus, an unknown language, and a
vast space of candidate readings let an analyst ``translate'' almost
anything, with internal consistency standing in for confirmation. The
canonical absurdity --- ``proving'' English is Semitic by cherry-picking
root matches --- recurs whenever the many comparisons a reading was
selected from go uncounted. The field's flagship survey enumerates
fifteen challenges, all properties of the object and its data; none
concern overfitting, multiple testing, false-positive control, or
contamination. We frame that methodological axis, modestly, as a
candidate sixteenth challenge and build machinery for it: a
falsification-first, decontaminated inference framework that treats
every positive as guilty until proven innocent. Four defences are
enforced --- a fabricated-language canary supplying a false-positive
floor, multiple-comparison deflation over signs, roots, and families
tried, a literature index separating published from literature-unseen
signs, and pre-registration with a held-out gate graded mechanically
rather than by the proposing model. This is a methods paper; we assert
no decipherment. Linear A is worked as the honest null: the binding
constraint is identifiability, not corpus size or effort. We report a
body of pre-registered nulls with their borderline cases named, and
analyse the June-2026 *301 ``cracked'' claim as a worked failure case.
The artifact is open and archived
(github.com/papadopouloskyriakos/logos; Zenodo concept DOI
10.5281/zenodo.21087572).

\hypertarget{introduction}{%
\subsection{1. Introduction}\label{introduction}}

Decipherment has a false-positive problem, and the field has no agreed
instrument for measuring it. The history of undeciphered-script research
is, in large part, a history of self-deception: a tiny corpus, an
unknown language, and a large space of candidate readings let a
determined analyst ``translate'' almost anything, with internal
consistency on the cherry-picked set standing in for confirmation. The
mechanism is not any single bad match but an unpaid multiplicity: with
enough candidate roots and a small enough corpus, a pile of
individually-plausible matches is the \emph{expected} outcome under pure
chance, so ``each individual comparison looks suggestive {[}while{]} the
conclusion is nonsense''. The canonical demonstration is deliberately
absurd --- one can ``prove'' that \emph{English is a Semitic language}
by matching a handful of English words to Proto-Semitic roots while
quietly ignoring the far larger number that do not match. The same move,
applied to a real undeciphered script, has produced a long line of
decipherment claims that no one can either confirm or refute.

That the discipline lacks a shared standard here is not our
editorializing but the field's own assessment. The flagship
computational-decipherment survey --- Braović et al.'s (2024) systematic
review of Bronze Age Aegean and Cypriot scripts --- enumerates fifteen
``major challenges,'' from unknown writing system and unknown language
to small dataset, incomplete vocabulary, allographs, and unavailable
parallel data (Braović et al.~2024). Those fifteen are, by construction,
properties of the \emph{object and its data} --- of the script, the
language, and the surviving corpus. What governs whether a \emph{claimed
reading} is believable is a different kind of thing: a property of the
inference methodology and its epistemics --- overfitting, multiple
testing, false-positive control, and contamination. We therefore frame
this axis, once and modestly, as a candidate \textbf{sixteenth
challenge} --- a deliberate rhetorical foil to the survey's fifteen, not
a claim to a new taxonomic slot or to sole discovery. It is
\emph{under-instrumented} rather than unrecognised: Packard's (1974)
nine deliberately fabricated decipherments already exercised exactly
this false-positive logic as a control, and the selective-inference
machinery we import (Bailey \& López de Prado 2014) is mature elsewhere
--- our contribution is to \emph{consolidate and operationalize} that
axis for decipherment, not to invent it. It is the methodological axis
the object-level taxonomy was never meant to cover, and it is distinct
from the survey's own acknowledged gap --- that on scoring results ``in
many cases there is nothing to compare them against, and their
evaluation remains an open problem'' (Braović et al.~2024, Sec. 6.3,
p.~746). That is an \emph{evaluation-target} problem (what to measure
against); the false-positive / multiplicity problem (whether a fitted
reading is being mistaken for a real one) is separable from it, and it
is the latter this paper instruments.

\textbf{Contribution.} We build the machinery for that axis: a
falsification-first, decontaminated inference framework for testing
undeciphered-script decipherment claims (Braović et al.~2024). The
harness treats every positive as guilty until proven innocent and
enforces four concrete defences: a fabricated-language canary that
supplies a hard false-positive floor; multiple-comparison deflation for
the number of signs, roots, and families a search was free to try; a
literature index that partitions published (\texttt{L\_known}) from
literature-unseen (\texttt{L\_virgin}) signs so that reproducing a known
value cannot masquerade as discovery; and pre-registration with a
held-out gate whose verdict is computed mechanically, never by the model
that proposed the reading. The evaluation gap gets a concrete, if
power-limited, answer: the recurring libation \emph{formula} --- roughly
50 inscribed objects spread across peak sanctuaries and caves at more
than five findspots --- supports a natural leave-one-site-out
cross-validation graded from held-out data. We are explicit that with so
few formula sites this cross-validation has low statistical power, and
we report it as such rather than as a decisive test (Braović et
al.~2024).

\textbf{This is a methods paper, and Linear A is its worked null.} We do
not claim a decipherment. If one is attainable on present evidence, the
framework is built to recognize it; when nothing survives, the
calibrated null --- what the corpus can and cannot support --- is itself
the deliverable. The working conclusion on current data is that the
binding constraint on Linear A is \textbf{identifiability}, not corpus
size or effort: there is no known related language to map onto, the
multiple-testing burden of candidate readings goes uncorrected, and the
sign inventory is mixed and partly lossy. A lossy, non-injective
sign-to-sound channel need not admit a decipherment at any corpus size,
so we make no claim that ``there is enough text to pin a reading''; the
rigorous demonstration of \emph{that} constraint is the result the field
can use. The June-2026 AI-assisted ``Linear A cracked'' claim ---
reading the \texttt{\textbackslash{}*301}-anchored libation verb as an
extinct Semitic form --- serves throughout as a worked example of
exactly these failure modes, analysed from its public, archived
presentation as a claim, not a person, and neither reproducible nor
held-out validated. The artifact is open and archived
(github.com/papadopouloskyriakos/logos; Zenodo concept DOI
10.5281/zenodo.21087572).

\textbf{Roadmap.} Section 2 situates logos against the Braović et
al.~taxonomy and the computational-decipherment lineage, including the
cognate-ID baseline. Section 3 specifies the harness --- canary,
deflation, literature index, and the pre-registration/held-out gate.
Section 4 develops the information-floor diagnostic and states its
identifiability limits. Section 5 works the June-2026
\texttt{\textbackslash{}*301} claim as a case study of the failure
modes. Sections 6--8 report the pre-registered nulls: morphology (§6);
metrology, phonology, and cross-script signals (§7); and contamination
and sufficiency probes (§8). Section 9 draws the discussion,
limitations, and conclusion.

\hypertarget{related-work}{%
\subsection{2. Related work}\label{related-work}}

\hypertarget{position-relative-to-the-survey-taxonomy}{%
\subsubsection{2.1 Position relative to the survey
taxonomy}\label{position-relative-to-the-survey-taxonomy}}

The reference frame for computational work on the Bronze Age Aegean
scripts is the systematic review of Braović, Krstinić, Štula and Ivanda
(2024, \emph{Computational Linguistics} 50(2):725--779), the field's
first survey focused specifically on Aegean and Cypriot scripts (its
only comparable predecessor being Ferrara and Tamburini 2022); the
broader machine-learning-for-ancient-languages landscape is surveyed by
Sommerschield et al.~(2023, \emph{Computational Linguistics} 49(3)).
Braović et al.'s Section 4 enumerates fifteen ``major challenges'' of
computational decipherment --- from unknown writing system, reading
direction and language, through small dataset, allographs and incomplete
vocabulary, to unavailable parallel data and hardware. None of the
fifteen concerns overfitting, multiple testing, false-positive control,
contamination, or the problem of distinguishing a \emph{fitted}
decipherment from a \emph{real} one. The same gap runs through Ferrara
and Tamburini's (2022, \emph{Lingue e Linguaggio} XXI.2:239--259)
predecessor survey: its §3.5 reviews the very cognate-matching lineage
logos builds on --- Snyder et al.~(2010) on Ugaritic↔Hebrew (the
29/30-sign benchmark our canary calibrates against), Berg-Kirkpatrick
and Klein (2011), and Luo et al.~(2019, 2021) --- and frames Linear A,
with Gelb and Whiting's typology, as a Type III script (unknown script
\emph{and} unknown language) whose ``necessary validation'' is
``altogether impossible,'' yet nowhere treats how a proposed match is to
be \emph{deflated} for the multiplicity it was selected from or
separated from a false positive. Two field surveys, the same unfilled
methodological axis. We do not score this against either as an
oversight: their challenges are properties of the object and its data
--- an unknown language, a small dataset, allographic variation, missing
parallel text --- whereas the axis logos foregrounds is a property of
inference \emph{methodology} and epistemics, a deliberately different
kind of thing. The false-positive / multiplicity axis is therefore
named-but-unfilled rather than absent, and logos supplies the machinery
to police it: preregistration, a decontaminated derivation/held-out
partition, multiplicity deflation, and a graduation gate computed
mechanically rather than by the model that proposed the reading.

Braović et al.~do gesture at a related but separable difficulty:
evaluation. Because computational analyses of an undeciphered script
have ``nothing to compare them against,'' they conclude (Sec. 6.3) that
``their evaluation remains an open problem.'' That is a distinct problem
from overfitting and multiplicity control, and logos addresses it
separately. Its held-out design is a concrete response --- the libation
formula recurs across a small set of peak-sanctuary and cave findspots,
permitting a leave-one-site-out cross-validation adjudicated
mechanically by \texttt{verdict.py} --- though we flag at the outset
that so few formula findspots give this particular test low statistical
power (quantified in Section 8). We locate our components in the
survey's own method taxonomy (Sec. 7.1) and, where prior art exists,
defer to it rather than claim novelty.

\hypertarget{the-current-philological-synthesis}{%
\subsubsection{2.2 The current philological
synthesis}\label{the-current-philological-synthesis}}

Salgarella (2025), \emph{Writing in Bronze Age Crete: ``Minoan'' Linear
A} (Cambridge Elements), by the author of the SigLA database, is the
most current authoritative one-author overview. It concludes that Minoan
is a language \textbf{isolate} --- it ``does not belong in any known
language family'' (the families used for comparison being Indo-European,
Semitic and Afro-Asiatic) --- typologically agglutinative with a
provisional Verb--Subject--Object order, and it records that ``no
bilingual text has so far been found'' in a corpus of roughly 2,500
inscriptions. Crucially, Salgarella ranks the etymological/comparative
method below combinatorial-contextual analysis and names
\textbf{digital/statistical analysis as ``the next desideratum in the
field.''} logos positions itself as exactly that desideratum --- but
under discipline, treating the isolate verdict as the binding
\emph{identifiability} constraint rather than a target to be argued
away.

\hypertarget{prior-computational-attempts-and-the-search-trap}{%
\subsubsection{2.3 Prior computational attempts and the search
trap}\label{prior-computational-attempts-and-the-search-trap}}

The seminal automatic-decipherment result is Luo, Cao and Barzilay
(2019), which recovers a character mapping from a lost script to a
\emph{known related} language via neural minimum-cost flow. It
auto-deciphered Ugaritic→Hebrew (+5.5\% over prior state of the art) and
translated 67.3\% of Linear B→Greek cognates, but did not address Linear
A and would fail there by construction: with no known cognate plaintext,
the objective has nothing to optimise toward. This cognate-language
precondition is the empirical form of our identifiability argument. The
cognate-ID baseline logos actually runs is \emph{not} a reimplementation
of Luo's neural flow: it is a published matcher, Tamburini (2025)'s
CSA\_OptMatcher (coupled simulated annealing), which reports the same
Luo 2019 accuracy metric. We adopt it as the non-neural baseline against
which any Linear A cognate claim must clear.

Against that disciplined baseline stand the in-domain instances of the
search trap: the dictionary-fishing programmes of Min Eu, Xu and
Cacciafoco (2019) and Loh and Perono Cacciafoco (2020), the latter
explicitly a ``brute force attack'' on dictionaries. These are the
concrete Linear-A cases the multiplicity correction is built to deflate.
On the visual side, cross-script comparison is \emph{not} our novelty:
Daggumati and Revesz (2019, 2023) already applied CNN+SVM cross-script
similarity, and Corazza et al.~(2022, ``Sign2Vec'') clustered signs
unsupervised; our defensible contribution is only the controlled
image-versus-sequence \emph{asymmetry} under one harness.

\hypertarget{the-discipline-lineage}{%
\subsubsection{2.4 The discipline
lineage}\label{the-discipline-lineage}}

Two older strands supply the honesty machinery. Packard (1974),
\emph{Minoan Linear A}, constructed nine deliberately fictitious
decipherments as a null control --- and the genuine Ventris-value
transfer beat them only about 2:1 --- anticipating our permutation-null
canary. The statistical lineage is selective inference and deflated
significance (Bailey and López de Prado 2014), the correction for the
number of hypotheses × signs × roots × families tried that turns a raw
match-rate into a defensible one; its transfer from the finance/backtest
setting to script decipherment is specified in Section 3. logos's empty
cell is the intersection none of these occupy: a falsification-first,
decontaminated inference framework that couples a published cognate-ID
matcher (Tamburini 2025, reporting the Luo 2019 metric) with
preregistration, multiplicity deflation and mechanical held-out
verdicts, released openly (github.com/papadopouloskyriakos/logos; Zenodo
concept DOI 10.5281/zenodo.21087572), with Linear A worked as the honest
null rather than a claimed crack.

\hypertarget{methods-the-discipline-harness}{%
\subsection{3. Methods: the discipline
harness}\label{methods-the-discipline-harness}}

The harness is a comparison-layer pipeline that a reader can
reimplement. A candidate reading enters as a \emph{proposal} and is
subjected, in fixed order, to five gates: pre-registration, deflation,
decontamination, held-out graduation, and a cross-cutting rule that
every gate is scored from persisted artifacts rather than from any
model's narration of its own success. Model-derived signals --- LLM
theses, learned embeddings, cognate-matcher assignments --- enter capped
at confidence ≤ 0.75, structurally unable to \emph{decide} a verdict;
only mechanically-verified held-out reads rank higher. The design
targets the false-positive / multiplicity / contamination axis, which
Braović et al.~(2024)'s fifteen-challenge taxonomy names but leaves
unfilled. Their fifteen challenges are properties of the object and its
data --- an unknown language, a small corpus, allographs; the axis logos
adds is a property of inference \emph{methodology} and epistemics,
deliberately a different kind of thing. (Their §6.3 note that
``evaluation remains an open problem'' gestures at evaluation, a
distinct problem from overfitting and multiplicity; the two are kept
separable here.)

\hypertarget{pre-registration}{%
\subsection{3.1 Pre-registration}\label{pre-registration}}

Before any model is built or run, the predictions are frozen. The
morphology pre-registration fixes, for each hypothesis, the claimed
affix, its source, the falsifiable test, and the acceptance threshold
--- with the git commit of the file serving as the timestamp, so the
model cannot be tuned to them after the fact. Each pre-registered affix
(drawn from Thomas 2020 and Duhoux/Valério) must recur on ≥2 distinct
stems across independent inscriptions, above a within-form permutation
null, and survive multiplicity correction at a target of \emph{p}
\textless{} 0.01 corrected. Outcomes are enumerated in advance ---
CONFIRM, EXTEND, REFUTE/NULL, or NO POWER --- so that a negative is a
publishable result, not a failed experiment. Any change is a new dated
pre-registration, never an edit.

\hypertarget{deflation}{%
\subsection{3.2 Deflation}\label{deflation}}

Internal fit on the derivation set is treated as no evidence. Every
match rate is deflated for everything tried:
\texttt{n\_hypotheses\ ×\ n\_signs\ ×\ n\_roots}. The nulls are
pipeline-level --- a frequency-banded permutation null (after Packard
1974) and an Euler-γ order statistic
(\texttt{expected\_max\_order\_stat(μ₀,\ σ₀,\ N)}) that gives the bar a
maximum over \emph{N} effectively-independent candidates should clear by
chance. Multiplicity is instrumented, not hand-passed: a deduplicating
\texttt{SearchLog} reports the exact distinct-candidate count
\texttt{n\_eff} a scanner produces (invariant 12). The strong structural
statistic, \texttt{S\_morph}, earns credit only through a productivity
gate --- an affix must attach to ≥2 distinct residual stems, so a single
formula word copied across inscriptions scores zero. After that fix, its
measured false-positive rate on within-form-shuffled corpora was
\textbf{3.0\% (6/200; Wilson 95\% CI {[}1.4, 6.4{]}\%)}, consistent with
the nominal one-sided \emph{z} ≥ 2 rate (≈ 2.3\%) within sampling error.
We flag this explicitly as an \textbf{indicative calibration check, not
the operative gate}: a 6/200 estimate is too coarse to certify an α =
0.01 threshold on its own, and it is reported only to show the detector
is not grossly mis-calibrated at the nominal level. The graduation gate
does not rest on it --- it applies the order-statistic bar over the far
larger permutation and candidate counts the deflation machinery tracks
(§3). Power was \textbf{100\% (50/50)} at a single strong planted-affix
strength --- a saturated point, not a full sensitivity curve (contrast
the phonology power curve, §7).

\hypertarget{decontamination}{%
\subsection{3.3 Decontamination}\label{decontamination}}

A frontier model has read Gordon, Best, and Di Mino, so a proposed
correspondence that merely reproduces a \emph{published} value is
shared-source contamination, not independent evidence. Three mechanisms
address this.

\textbf{The literature index} records published Linear A sign-value and
correspondence claims; any proposal matching the index is tagged
\texttt{literature\_match} and quarantined --- it may be tested but can
never count as discovery. The index partitions the sign inventory into
\texttt{L\_known} (referenced by ≥1 published proposal) and
\texttt{L\_virgin} (untouched). Discovery claims may rest only on
\texttt{L\_virgin} held-out success, since regurgitation can only return
what is in the literature. The seed is deliberately conservative --- a
sign wrongly omitted is wrongly granted \texttt{L\_virgin} status, the
dangerous direction --- and every disputed entry is indexed to
\emph{catch} regurgitation, never as an accepted reading. Population is
adversarial: of the seeded West-Semitic lexical proposals, five are
Gordon's equations and a-sa-sa-ra-me = ``oh Asherah!'' is Best (1981),
not Gordon; candidates failing independent corroboration against
Rendsburg (1996) (\texttt{qa-pa} = kappu, a Semitic \texttt{ki-ro}) were
rejected rather than indexed. Generated from the seed (invariant 12),
the index currently holds \textbf{63 published claims across 62 indexed
entries (55 single signs)} --- a deliberately conservative,
non-exhaustive seed, not yet the full §C.1 literature census.

\textbf{The L\_fake canary} is a phonotactically-plausible invented
lexicon, calibrated to a real candidate language's phoneme frequencies,
root-template structure, and lexicon size, but never published and never
in any training corpus. Running the full pipeline against it yields an
empirical false-positive floor: any signal it produces is definitionally
spurious, and memorisation cannot be invoked because there is nothing to
memorise. A real candidate's match must clear the L\_fake score
distribution by the corrected margin. Its own divergence from the
calibration targets is reported, never minimised; with a Semitic
root-template sampler and an external Hebrew (ETCBC/bhsa) reject set,
the root-template total-variation distance fell from \textasciitilde0.84
to \textasciitilde0.33 end-to-end (\textasciitilde0.23 standalone) on
the Ugaritic↔Hebrew gold, and any residual real-lexicon collision is
measured and published rather than asserted to be zero.

\textbf{LLM-ablation} compares the pipeline with and without the model's
proposals to isolate contamination --- the remaining gate on the
contamination number, matching a model's cognate glosses against the
indexed lexical claims.

\hypertarget{the-held-out-graduation-gate}{%
\subsection{3.4 The held-out graduation
gate}\label{the-held-out-graduation-gate}}

A reading graduates only mechanically, from held-out data, as a
conjunction of pre-registered clauses computed by
\texttt{scripts/verdict.py} --- never by the model that proposed the
reading. The operative deflation clause is an \textbf{order-statistic
bar}: the observed held-out statistic must exceed the \emph{expected
maximum} of the multiplicity-null taken over the number of trials the
search was free to run (the E{[}max{]} over \texttt{n\_trials}
deflation). This is a direct correction for selection across the
hypothesis space, and it fails closed on a degenerate null (zero
variance ⇒ no graduation). The Deflated-Sharpe Ratio (Bailey \& López de
Prado 2014) --- a probability in {[}0, 1{]}, the confidence that the
true deflated performance exceeds its multiplicity-adjusted benchmark
--- is computed and \textbf{reported} alongside; after pre-submission
review we made the order-statistic bar the \emph{operative} clause
rather than a fixed DSR ≥ 0.95 cutoff, because the expected-max bar
needs no normal approximation and fails closed on degenerate input. The
match-rate the deflation is applied to is supplied by a cognate-ID
baseline that reports the Luo et al.~(2019) accuracy metric; the
instance we run is Tamburini's CSA\_OptMatcher (coupled simulated
annealing, §2.3--2.4, reproduced in the artifact), so the gate does not
depend on any single citation. A second pre-registered clause,
\textbf{\texttt{generalizes\_to\_virgin}}, requires the reading to
recover literature-unseen (\texttt{L\_virgin}) signs above a
pre-committed threshold and on a minimum number of distinct such signs
--- so reproducing already-published values cannot by itself graduate a
claim, and the clause fails closed when no threshold or no supported
virgin sign is present. The L\_fake corpus must \emph{fail} as a null,
its corrected margin must be cleared, and the observed statistic must
beat the order-statistic bar over the Packard/random-lexeme multiplicity
null (the within-inscription sign-order shuffle is a null in the
morphology module, §6, not a gate clause). The
minimum-description-length check (k ≤ U\_floor, §4) was, after the same
review, \textbf{removed from the gate and is now reported as a
diagnostic rather than enforced as a graduation clause}: clearing it
does not certify identifiability (it is a necessary-precondition /
toy-model diagnostic, §4, §9), and as a hard clause it could only be
evaluated against a free-parameter default that risked passing
degenerate inputs. U\_floor is still displayed next to every claim, and
``parameters \textgreater{} information ⇒ underdetermined'' is still
stated (invariant 7) --- it simply no longer sits inside the AND. These
are \emph{gate parameters} --- pre-registered choices for this corpus,
not universal constants. Validation is leave-one-site-out, not ordinary
\emph{k}-fold, because the libation formula induces formulaic dependence
across inscriptions. The gate is AND-monotone --- clauses can only make
it stricter --- so hardening never lets a previously-failing reading
graduate.

\hypertarget{grade-from-persisted-artifacts}{%
\subsection{3.5 Grade from persisted
artifacts}\label{grade-from-persisted-artifacts}}

The invariant across every gate: verdicts are computed from persisted
rows, never from a model's account of its own outcome. The verdict
writer is grep-clean of any model call and is the sole writer of the
verdicts table. This discipline caught confounds inside logos's own
builds --- a \texttt{S\_morph} statistic that conflated repetition with
morphology, a verdict clause that was a tautology and could never fire
False, and a search-log tripwire that flagged honest multi-sign logs.
Each was fixed and locked by regression tests before commit,
illustrating the harness turned on itself. The artifact is public and
archived (repository https://github.com/papadopouloskyriakos/logos;
Zenodo concept DOI 10.5281/zenodo.21087572).

\hypertarget{the-information-floor}{%
\subsection{4. The information floor}\label{the-information-floor}}

Before any decipherment method is graded, the corpus must be sized
against what it can identify. Two quantities matter: how much text
exists, and how many free parameters a reading spends against it.
Confusing the first for the second is the origin of the ``too little
text'' folk claim --- and, we argue, of its equally unearned inverse,
the optimism that a big-enough corpus makes a reading identifiable.

\textbf{The corpus.} The current authoritative synthesis (Salgarella
2025) gives an envelope of roughly 2,500 total Linear A finds, of which
about 1,400 are readable inscriptions, drawn from 56 findplaces spanning
ca. 1800--1450 BCE (Middle Minoan II--Late Minoan IB) but overwhelmingly
concentrated in the single LM IB destruction horizon (Salgarella 2025
§1--3, §6). That near-single horizon is why a site-held-out split, not a
chronological one, is the only realistic partition. No bilingual text
has ever been found (Salgarella 2025 §8).

\textbf{Which denominator?} The recurring sign-count confusion resolves
into four distinct denominators, each of which must be labelled where it
is used (Salgarella 2025 §1--3): \textbf{97} (Braović/Tan token-frequent
syllabograms), \textbf{\textasciitilde150} (full syllabary estimate),
\textbf{259} (the logos working repertoire, conflating syllabograms,
subscripts, and logograms), and \textbf{\textasciitilde350} (total
including ideograms; GORILA catalogues 390 signs). The channel through
which structure is observed is a \emph{sign-occurrence} count, not a
document count or a sign inventory (invariant 7), and two occurrence
counts must be kept distinct. The broader \textbf{\textasciitilde7,400
sign occurrences} across the whole corpus (Schoep 2002; Salgarella 2025
§1--3) counts every sign type --- syllabograms, logograms, and numerals
alike. The narrower \textbf{N ≈ 5,792 parsed syllable-signs} is the
subset that the unicity computation below actually consumes, after
logograms and numerals are stripped. The two are different channel sizes
and are used in different places; the \textasciitilde7,400 figure never
enters the unicity arithmetic. Two further properties tighten the
budget: Linear A words are short and formulaic (the six-position
libation-formula template; Salgarella 2025 §7), and a single
language/horizon dominates.

\textbf{The unicity argument, and its limit.} On the logos-normalized
corpus (1,341 inscriptions, N ≈ 5,792 parsed syllable-signs, V = 259),
Shannon's unicity-distance criterion yields a measured range of
\textbf{U ≈ 204--415 signs} across the whole plausible (V, P) parameter
space, everywhere far below N. \textbf{We retain this result only as a
toy-model precondition diagnostic, and explicitly withdraw it as any
sufficiency or refutation claim.} It does \emph{not} show there is
enough text to pin a decipherment, and we do not argue that. Unicity
distance presupposes a \emph{known plaintext-language oracle}; because
we estimate the redundancy D from the ciphertext itself, that oracle
collapses into Linear A up to relabelling, so the computation
demonstrates only measurable recurrent structure, not that any reading
is identifiable. The precondition it checks is necessary, not
sufficient. And it can fail to exist entirely: if Linear A underspells
like Linear B, the sign→sound channel is lossy and non-injective, and
such a channel may possess no unicity distance at all. Corpus size, in
short, cannot settle the question either way.

\textbf{What actually binds.} The binding constraint is
\emph{identifiability}, not corpus size (Salgarella 2025 §8). It has
three parts: (i) the \textbf{absence of a known cognate language} ---
the map has a value-space but no anchored target (Luo 2019's limit);
(ii) the \textbf{multiple-testing / search trap} --- that a unique
consistent map may exist in principle neither guarantees an algorithm
finds it nor stops a cherry-picked false positive across the parameters
tried (invariant 8); and (iii) a \textbf{mixed inventory} of
syllabograms, logograms, and numerals that a real reading must first
separate. These, not the count of surviving signs, are what a
decipherment must overcome. Accordingly, when a candidate lexicon grows
without bound on a fixed occurrence budget --- free parameters exceeding
the floor --- the reading is underdetermined, and the framework must say
so.

\hypertarget{a-worked-failure-case-the-301-claim}{%
\subsection{5. A worked failure case: the *301
claim}\label{a-worked-failure-case-the-301-claim}}

The framework of Section 3 is abstract until it meets a concrete claim.
In June 2026 a widely circulated ``Linear A cracked'' announcement (Di
Mino 2026) argued that Linear A records an extinct Semitic language,
reviving Gordon's (1957, 1966) hypothesis. Its load-bearing anchor is a
single reading: the recurring peak-sanctuary libation formula word
A-TA-I-*301-WA-JA, in which the Linear-A-only sign *301 is \emph{itself}
assigned a /na/-initial value, so that the verb root reads as
West-Semitic N-W-Y ``to dwell'' (\emph{nawaya}) and is thereafter
matched to later Hebrew liturgy. We analyse this \textbf{as a claim, on
its public evidence} --- not the person --- from its public and archived
presentation; our reference copy of the claim and its table is the
archived record carried in the logos artifact (v0.1.0), and the
repository (github.com/papadopouloskyriakos/logos; Zenodo
10.5281/zenodo.21087572) carries the machinery that indexes it. It is a
clean, self-contained instance of exactly the failure modes the gate is
built to catch. The refutation is three-part and rests, wherever
possible, on the claim's \textbf{own published table} and on primary
sources. We flag the meta-point explicitly: this refutation is a
literature-and-primary-source argument, not itself a logos held-out
result --- it critiques the claim on shared evidentiary ground rather
than adjudicating it mechanically through our own gate. That gap is the
sharpest limitation of the present paper, and we name it: the graduation
gate and the L\_fake canary are validated as instruments (on unit tests
and on the known-answer Ugaritic↔Hebrew surrogate) but have not been
\emph{run on this claim}. The natural demonstration --- scoring Gordon's
West-Semitic equations and the Di Mino lexicon against the L\_fake
false-positive floor and the order-statistic bar, so the reading either
clears the machinery or fails closed against it --- is directly runnable
with the components of §3 and is the priority of the planned revision
(docs/findings/2026-07-01-referee-simulation). Until then,
``falsification-first'' is demonstrated on synthetic positives and
argued (not yet mechanically adjudicated) on this real one.

\textbf{(a) The segmentation is not novel.} Root i-*301 plus affixation
is the reading of Davis (2013); an independent segmentation by Thomas
(2020) reaches the same core and is tested head-to-head against it (the
current synthesis, Salgarella 2025, does not itself adopt Thomas). The
affix-stacking is proven by multi-attestation forms ---
ta-na-i-*301-u-ti-nu (IO Za 6) and a-ta-i-*301-de-ka (ZA Zb 3) --- that
vary the material around a stable i-*301 core. (These sigla are
confirmed against Davis 2013: \texttt{TL\ Za\ 1}, \texttt{IO\ Za\ 6},
and \texttt{ZA\ Zb\ 3}.) The claim re-segments a structure already
established in the literature. Under the decontamination layer this
matters procedurally, not just historically: a correspondence that
reproduces a published segmentation is shared-source contamination,
tagged \texttt{literature\_match} and barred from counting as discovery
(logos v0.1.0, §C.1).

\textbf{(b) The gloss is contested and weaker.} The mainstream value of
the verb is ``give/dedicate'' --- Duhoux's (1992) gloss, adopted by
Davis (2013) as a \emph{structural} reading that assigns *301 \textbf{no
phonetic value at all}, and endorsed by the current synthesis
(Salgarella 2025, §7.2). ``Dwell'' is a semantically distant,
unmotivated substitution. The point is not that ``give'' is certainly
right, but that a \emph{structural} reading needing no sound-value for
*301 already accounts for the data the phonetic claim is invoked to
explain.

\textbf{(c) The primary descriptions of this form's morphology do not
support the Semitic label.} The claim's own table segments the verb as
prefix-stacking and agglutinative --- {[}prefix A ``I''{]}
{[}stem-morpheme TA{]} {[}stem-vowel I{]} {[}root *301{]}. We
deliberately do \textbf{not} rest on an \emph{a-priori} typological
contradiction here, and we retract the earlier draft's claim that a
prefix-stacked verb is ``internally incoherent'' with a Semitic reading:
Semitic verbal morphology includes productive prefix conjugations, so a
prefix on the verb is not by itself un-Semitic. The point is empirical,
not typological, and it comes from the primary literature. Davis (2013,
p.~38 n.~13) characterises i-*301's morphology as ``neither IE nor
Afroasiatic'' --- and Semitic \emph{is} Afroasiatic, so a decade earlier
the primary description already found \textbf{no Afroasiatic
morphological signature} in this very form. (We rest on what the
footnote states --- the absence of an Afroasiatic signature --- not on a
stronger ``ruled out'' than a hedged footnote can bear.) Independently,
Salgarella (2025, §8) concludes Minoan is agglutinative (after Duhoux
1978) and a language isolate. The claim's own segmentation is thus
consistent with the isolate typology, while the specific Semitic
assignment finds no morphological support in the primary descriptions of
the form.

\textbf{The free-parameter receipt.} The decisive point for
identifiability is that *301 has no cross-script anchor: in the claim's
own table its AB-number column is ``---'' and its Linear B
correspondence ``(none)''. Even the standard practice of reading Linear
A homomorphs with their Linear B values is a backward projection to be
handled with care (Steele \& Meissner 2017; already cautioned by Duhoux
1989, apud Davis 2013, p.~38 n.~11) --- and an A-only sign such as *301
has no homomorph to project from at all, so nothing cross-script
constrains its value. The single sign on which the entire decipherment
pivots is thus a \textbf{free parameter} --- its value was assigned to
fit the desired root, not constrained by any independent evidence. This
is precisely the case the minimum-description-length / identifiability
\emph{diagnostic} flags: a value chosen to make one word read out
carries no description-length saving that survives the multiplicity of
alternatives it was selected from (cf.~Section 3; the diagnostic is
reported next to every claim, and the operative graduation clause ---
the order-statistic bar over exactly that multiplicity --- is one a
hand-fitted value cannot clear). The binding obstacle is not corpus size
but identifiability --- no known cognate language to constrain the
assignment, uncorrected multiplicity, and a sign whose value is chosen
rather than anchored.

Read against Section 3, the claim instantiates two distinct failure
modes at once. It is \textbf{narrative capture} --- the table's
\emph{structure} says agglutinative and prefixing, but the \emph{label}
says Semitic, and the label wins because the desired conclusion (a
Semitic prayer) is a fixed prior rather than an emergent reading. And it
is a \textbf{free-parameter / identifiability} failure --- an unanchored
sign-value tuned to the target, on a corpus with no known cognate
language to constrain it. logos indexes the reading (*301 = /na/)
precisely so a model cannot regurgitate it as discovery, and quarantines
rather than dismisses it. The lesson generalises: a match rate on a
hand-picked word, absent a committed prediction and multiplicity
accounting, is not evidence. The *301 claim is the worked example of why
the gate must fail closed.

\hypertarget{results-i-morphology-induction}{%
\subsection{6. Results I: morphology
induction}\label{results-i-morphology-induction}}

The first worked application of the framework is pre-registered,
unsupervised morphology and segmentation induction over the full Linear
A corpus (1,341 inscriptions across 52 sites --- the 52 that survive
logos normalization of Salgarella's 56-findplace envelope). The target
claims are the recurring affixes and positional grammar of the
Davis/Thomas/Duhoux--Valério tradition: prefixes and suffixes
(\texttt{i-}, \texttt{a-}, \texttt{-te/-ti}, and Davis's
\texttt{i-\textbackslash{}*301} verb stem) that field analyses read as
genuine Minoan morphology. The pre-registration was frozen before the
model touched the data, and every candidate was graded mechanically
against two null floors: a within-word sign-order permutation (shuffle)
floor, and an \texttt{L\_fake} floor built from a first-order
sign-bigram (Markov) corpus with zero lexical overlap with the real
inscriptions. The design earns the paper's core discipline --- internal
recurrence is not evidence; only separation from a generative null that
reproduces sign statistics counts.

\textbf{Pooled run: an honest null.} Of the 16 pre-registered affixes, 4
beat the within-form permutation null and the deflated multiplicity bar
(\texttt{i-\textbackslash{}*301}, \texttt{a-}, \texttt{i-},
\texttt{-te/-ti}), a real confirm rate of 0.250 (Wilson 95\% CI {[}0.10,
0.50{]}). This clears the shuffle floor decisively: there \emph{is}
recurring edge structure, exactly as Davis/Thomas describe. The apparent
standout is \texttt{i-\textbackslash{}*301}, which sits far above the
permutation null (six distinct affix environments, z ≈ 8.0,
DSR-corrected p ≈ 1e-14) --- \textbf{but that figure is computed against
the weaker within-form permutation null only}, and it does not survive
the operative \texttt{L\_fake} bigram floor. Under the bigram floor the
whole set fails: the fabricated Markov corpus confirms the same affixes
at rate 0.375 (Wilson 95\% CI {[}0.18, 0.61{]}), above the real 0.250.
On this corpus the apparent affixation ---
\texttt{i-\textbackslash{}*301} included --- is fully reproducible by a
sign-transition model, so \texttt{has\_morphology\_power} is False and
the module reports the null. This is a limit of the \emph{data}, not of
the field's analysis: Linear A words are short --- mean 1.84 signs over
all word-\emph{tokens}, median and mode 1, with 76.1\% of word-tokens at
≤2 signs (56.5\% single-sign) --- and on such words ``morphology'' and
``bigram order'' are statistically indistinguishable. The premise
reconciles with Fuls's 3.3-sign figure once the denominator is matched:
Fuls (2015:57, §6.5) computes the mean over a word \emph{list} of ``554
complete sign sequences, where duplicates are reduced to one'' (GORILA;
Godart and Olivier 1985) --- i.e.~over distinct word \emph{types}, not
tokens. Restricting our corpus to distinct multi-sign types likewise
recovers 3.07; the token vs.~type denominator is the whole discrepancy.
We note one tension worth flagging rather than burying: Fuls reads the
3.3-sign length, via his exponent-to-word-length regression, as evidence
that Minoan is \emph{polysynthetic} (rich affixation); our direct
morphology test finds the apparent affixation bigram-reproducible with
no power, so it does \textbf{not} corroborate a rich-morphology reading
--- the word-length typology and the direct null point opposite ways,
and we report the null. This bigram free-rider is one of the would-be
confounds the machinery has caught and reported as a null rather than a
finding. One validation this null still lacks --- and the first item of
the planned revision --- is a \textbf{positive control}: metrology and
phonology (§7.1--7.2) each demonstrate their detector fires on a planted
signal, but the morphology bigram-floor test has not yet been shown to
recover \emph{known} affixation, so ``no power on Linear A'' cannot yet
be cleanly separated from ``the floor is too strong.'' The direct
control is the identical bigram-floor test on \textbf{Mycenaean Greek
(Linear B)}, whose longer words carry uncontested inflectional
morphology; the DĀMOS Linear-B corpus ingested for §7.3 (13,562
wordforms) now makes it runnable. Our short-word argument makes a
falsifiable prediction there --- the test \emph{should} regain power on
Linear B --- which, if borne out, would establish the Linear A result as
a short-word / low-power situation rather than a claim that Minoan lacks
morphology; until that control is run we report the Linear A outcome
strictly as \emph{no power}, not as evidence against affixation.

\textbf{A modest segmentation positive.} The word-boundary segmenter is
a genuine, held-out positive. Under leave-one-site-out cross-validation
(52 sites; not k-fold, given formulaic dependence), a DP-unigram
(Goldwater-style) segmenter recovers the scribe's own word divisions at
micro-F1 0.436 versus a random-boundary baseline of 0.389 (boundary base
rate 0.406). Morfessor over-segments the short syllabic corpus (micro-F1
0.156) and is reported as not usable rather than hidden. The
\emph{method} works above chance even where the \emph{morphology claim}
does not survive the bigram floor.

\textbf{Stratified re-run.} Because pooling the corpus risks inducing
``morphology'' that is really a mix of genres --- a confound the
pooling-preserving permutation null cannot catch --- a second
pre-registered addendum re-ran induction per genre stratum. The strata
reconcile exactly to the corpus total: 290 admin + 99 libation + 129
other + 740 seal (structurally excluded as an abbreviation channel) + 83
emptied by abbreviation-strip = 1,341 inscriptions. Induction ran on
admin (290 inscriptions, 1,486 words, 14 sites), libation (99, 370, 15),
and other (129, 187, 40). Note that the 99-inscription ``libation''
\textbf{genre} stratum --- stone-vessel inscriptions over 15 sites ---
is \emph{not} the same object as the libation-\textbf{formula} corpus
(\textasciitilde50 inscribed objects drawn from peak sanctuaries
\emph{and} caves, across more than five findspots) that carries the
recurring dedicatory formula; conflating the two overstates the
evidence. Following Salgarella \& Judson (2024), chronological
cross-validation was explicitly declined (988/1,341 ≈ 74\%
single-horizon LM IB), and site --- not scribal hand, which is not
individuated for Linear A (Salgarella 2019) --- is the independence
unit, gated by presence at ≥2 distinct sites. The candidates are tested
for cross-stratum stability above a \emph{per-affix} \texttt{L\_fake}
floor. The result is \texttt{has\_validated\_morphology\ =\ False}: no
pre-registered affix is cross-stratum stable above the bigram floor at
the pre-registered α = 0.01. Stratification also showed that even
\texttt{i-\textbackslash{}*301} is bigram-reproducible in the repetitive
libation register, where the Markov corpus manufactures the apparent
verb morphology; Davis's slot grammar is not refuted, merely not
separable from sign bigrams here. Where formula-bearing sites do drive a
signal, statistical power is intrinsically low: a leave-one-site-out CV
over so few formula findspots (an effective \textasciitilde5-fold site
partition) cannot resolve a modest effect, and we say so rather than
read the borderline result as support.

\textbf{Sensitivity, disclosed.} An independent adversarial review
corrected an initial overclaim of robustness. The null is a threshold
call, and the honest 2×2 is: under both α thresholds at the defensible
bucketing (libation = stone vessels only), NULL; under loose bucketing
(libation including stone objects and inked), NULL at α = 0.01 but at α
= 0.05 the affixes \texttt{i-} (locative) and \texttt{-te/-ti}
(ablative) validate. These are precisely the Duhoux/Valério H3 nominal
affixes that Salgarella (2025) §8 independently endorses --- the named
borderline candidates, the closest thing to recoverable morphology in
the corpus, but fragile to \emph{both} the α threshold and a defensible
genre choice. The pre-registered α = 0.01 null is itself robust to the
bucketing, so the bucketing fix is not cherry-picking the null.
Separately, logos's own self-review --- not the adversarial reviewer ---
caught the bigram free-rider \texttt{a-} and excluded it via the
per-affix floor; the adversarial reviewer also corrected a docstring
that mislabelled the ≥2-site presence gate as effective-n deflation. We
report the null; full effective-n deflation of the z-test --- which
would only strengthen it --- is a recorded, deferred limitation.

\hypertarget{results-ii-metrology-phonology-and-the-cross-script-asymmetry}{%
\subsection{7. Results II: metrology, phonology, and the cross-script
asymmetry}\label{results-ii-metrology-phonology-and-the-cross-script-asymmetry}}

Three further pre-registered probes exercise the harness against
distinct field claims. Each returns a calibrated null or a truth-capped
positive; together they map where the Linear A corpus does and does not
carry recoverable signal, and they locate the binding constraint at
identifiability rather than effort.

\hypertarget{metrology-a-parse-coverage-null-that-agrees-with-the-field}{%
\subsubsection{7.1 Metrology: a parse-coverage null that agrees with the
field}\label{metrology-a-parse-coverage-null-that-agrees-with-the-field}}

Direction D poses the accounting tablets as a constraint-optimisation.
We parse HT (LM I) documents into balance units of the form
\texttt{{[}recipient{]}{[}commodity\ logogram{]}{[}integer{]}{[}fraction\ signs{]}}
and solve the fraction-sign rational values by RANSAC over exact
\texttt{Fraction} Gaussian elimination --- robust to the
unbalanced-tablet outliers that defeat least-squares --- requiring line
items to balance the preserved KU-RO totals. Validation is grouped
k-fold by document (no tablet straddles train and test) against a
permutation null.

On the real corpus (32 documents / 35 balance units), integer-level
balance holds for 7/28 (25\%). Of seven fraction-bearing constraints,
only one is satisfied by the solved system, and the sole value our
balance method can determine is the fraction sign J = 1/2, forced by
HT104's \texttt{2·J\ ==\ 1}. Held-out fraction balance is 0.0 against a
permutation-null mean of 0.113 (max 0.143), p = 1.0, not separated: the
solved fractions do not generalise above the shuffle null. A synthetic
positive control confirms this is not a machinery bug --- with planted
values the solver recovers all of them and reaches held-out balance 0.9
vs null 0.22, p = 0.003, separated. The null is driven by sparse
automated parse coverage: only \textasciitilde7 fraction-bearing
constraints survive automated parsing of 32 tablets, because
multi-commodity blocks, illegible number glyphs, KI-RO deficit
sub-lists, and editorial restorations are not captured.

Crucially, this is a null the field itself reached. Corazza et
al.~(2021, p.2) abandoned the balance-to-totals method precisely because
no Linear A inscription containing numeral-phrase totals is free of
reading or calculation problems, moving instead to constraint
programming over ordinal/typological premises. logos's parse-coverage
null is therefore consistent with the published experts, not a
logos-unique failure. The claim is method-specific: our balance method
determines only J = 1/2 --- but that value is not novel, it is the
long-standing consensus (Bennett 1950), corroborated here independently
by Corazza et al.~(2021) and by Schrijver (2014). Schrijver's three
further J = 1/2 tablets (HT Zd 156, PE 1, HT 123a) lift the value to at
least four corroborating documents. The result illustrates the harness
working as designed: a claim that fails held-out validation is reported
as a null, and the one value that survives is credited to its actual
originator rather than re-branded as a logos discovery.

\hypertarget{distributional-phonology-a-data-limited-null-bounded-by-a-power-curve}{%
\subsubsection{7.2 Distributional phonology: a data-limited null bounded
by a power
curve}\label{distributional-phonology-a-data-limited-null-bounded-by-a-power-curve}}

Track 2 asks whether a sign's distributional context --- which signs sit
beside it --- carries information about its phonetic value, measured on
the known \texttt{AB} signs whose Linear B sound values can be borrowed.
Features are PPMI-weighted, L2-normalised left/right co-occurrence
counts; labels are the Linear-B-transferred (C, V) parsed from the
GORILA name, so features never see the labels. The test is leave-one-out
1-NN against a label-permutation null, Šidák-deflated over the two
tests; 50 of \textasciitilde55 AB signs met the ≥3-occurrence threshold.

Neither test clears chance: consonant class (13 classes) at 0.080
accuracy (null 0.065, permutation p = 0.40) and vowel class (5 classes)
at 0.140 (null 0.201, permutation p = 0.86). The honest verdict is
\emph{data-limited}, not \emph{no bridge exists}, because the positive
control disciplines the reading. A first draft's flat positive control
failed to fire at planted strength 1.0; replacing it with a power curve
over planted strengths shows the test firing (p \textless{} .05) at
strength ≥ 2 but not at 1. The test works but is not infinitely
sensitive at n = 50 / 13 classes, so we cannot distinguish ``no
distribution→phonetics bridge'' from ``too few signs to detect a weak
one.'' Both point the same way: the sound path is not recoverable from
Linear A alone. The obvious next move --- impute sound values
cross-script from a larger Linear B corpus --- is now \emph{tested
rather than deferred}: §7.3 reports that ingesting the full DĀMOS
Mycenaean corpus leaves the cross-script distributional alignment at
chance, so more Linear B data is necessary but not sufficient; the block
is an absent distributional signal, not corpus volume alone. No sound
value is imputed for any sign. This pre-empts the predictable ``couldn't
you just use ML to recover the sounds?'' objection with a quantified
detectability floor rather than hand-waving.

\hypertarget{the-cross-script-asymmetry-image-alignment-recovers-shared-anchors-but-the-recovery-is-circular-by-construction}{%
\subsubsection{7.3 The cross-script asymmetry: image alignment recovers
shared anchors, but the recovery is circular by
construction}\label{the-cross-script-asymmetry-image-alignment-recovers-shared-anchors-but-the-recovery-is-circular-by-construction}}

The palaeographic probe renders 88 Linear-A and 74 Linear-B glyphs and
tests held-out recovery of shared A↔B anchors. Classical descriptors
(HOG + Hu moments + shape-context) reach direct-NN accuracy 0.410
{[}0.18, 0.64{]} and Procrustes 0.185 {[}0.00, 0.36{]}, far above image
chance (0.0135); a from-scratch I-JEPA reaches only 0.324 ± 0.025 (3
seeds) and is excluded from the claims as over-parameterized for
\textasciitilde90 glyphs. Image embeddings recover held-out anchors
where sequence co-occurrence sat at chance. Two controls decide how to
read that asymmetry, and together they \emph{downgrade} it from a
positive to a demonstration.

\emph{First, a font-swap control.} Rendering both scripts in one
typeface (Aegean, George Douros) risks making the alignment a
single-designer house-style artifact rather than a fact about the signs.
We re-render Linear A in Noto Sans Linear A and Linear B in Noto Sans
Linear B --- two independent type designers, each font covering only its
own Unicode block (Noto-A carries zero Linear B glyphs and vice-versa),
so no single hand draws both scripts. The recovery survives and is if
anything stronger (direct-NN 0.367→0.416, Procrustes 0.156→0.195 on the
control's matched anchor set; both \(\gg\) chance): the font-artifact
hypothesis is \emph{refuted}, and ``distances reflect shape, not font''
is now demonstrated rather than asserted
(\texttt{scripts/palaeo/font\_control.py}).

\emph{Second, and decisively, the circularity the font control cannot
remove.} Linear B was historically adapted from Linear A, and the A↔B
transcription \emph{values} that define the anchor set were assigned by
twentieth-century epigraphers largely on the basis of sign-shape
similarity. An image encoder that ``recovers the A→B value map''
therefore re-derives a correspondence that was \emph{defined} by shape
--- the result is close to tautological and is not independent evidence
for any phonetic reading. We accordingly relabel this leg a
\textbf{shape-genealogy engineering demonstration, not an offensive
positive}: it confirms a known palaeographic fact (borrowed signs look
alike) and validates a font-robust descriptor pipeline, but it reads no
Minoan. It remains truth-layer-capped in any case --- the Ferrara,
Montecchi \& Valério ``Archanes Formula'' hard negative
(\texttt{CH\ 095} resembles \texttt{AB\ 60}/\emph{ra} but is not its
phonetic counterpart) shows a plausible shape match can be phonetically
wrong, so a shape match is a ≤ 0.75 signal (invariant 5), never a
reading. The honest asymmetry is thus: the image path re-derives a
shape-defined mapping (expected, circular, capped), while the sequence
and cognate paths \textbf{remain at chance even after ingesting the full
DĀMOS Mycenaean corpus} --- 13,562 wordforms, ≈15× the earlier cognate
file, 56 anchors: the best of five alignment methods reaches 0.025
against a 0.011 chance floor
(\texttt{clearly\_above\_chance\ =\ False}), while the \emph{positive
control recovers 0.947}, so the harness works and the null is a real
null. Ingesting all of DĀMOS therefore reframes the obstacle from
\emph{data volume} to \emph{absent distributional signal}: the A↔B
phonetic correspondence is simply not carried by cross-script sign
co-occurrence, even at full Linear-B scale. And the A-only signs on
which offensive claims would pivot (e.g.~*301) have no cross-script
anchor at all, so images cannot impute them even in principle.
\textbf{The study reports no offensive positive.}

\hypertarget{results-iii-contamination-and-information-sufficiency-probes}{%
\subsection{8. Results III: contamination and information-sufficiency
probes}\label{results-iii-contamination-and-information-sufficiency-probes}}

Two failure modes are invisible to the false-positive / multiplicity
axis of §7: a decipherment that \emph{looks} novel but silently
reproduces published readings a language model has memorised, and a
corpus that may be too small --- or too unidentifiable --- to recover
any reading at all. This section reports the decontamination and
information-sufficiency probes that target them. Both are run as method
demonstrations. The LLM-ablation and \texttt{L\_virgin} probes return
completed (if deliberately bounded) findings; the
information-sufficiency sweep now returns a \textbf{measured
accuracy-vs-size curve} on a corrected matcher, reported here with its
larger-size points still completing (§8.3, marked \emph{under
revision}). Throughout, where a clean number was not yet obtainable we
report the partial-null and the confound rather than a headline.

\hypertarget{llm-ablation-does-a-frontier-model-regurgitate-the-literature}{%
\subsubsection{8.1 LLM-ablation: does a frontier model regurgitate the
literature?}\label{llm-ablation-does-a-frontier-model-regurgitate-the-literature}}

We ran the §C.4 ablation against \texttt{qwen2.5:72b} --- a
frontier-scale open model, a far stronger memoriser than a local 12B ---
on a rented H100 (40 seeds × 40 held-out Linear A forms; NW-Semitic
candidate; a Hebrew lexicon of 7,682 skeletons for the model-free arm; 0
LLM errors across 40/40 seeds) (logos v0.1.0). The first-pass
contamination rate came out at \texttt{1.000} in all 33 seeds where it
was defined --- but this is an \textbf{artifact, reported as such}: the
mechanical arm never shares a literature-matching correspondence with
the LLM (their value vocabularies overlap in only 10 of 134 values), so
the rate is forced to 1.0 by construction. Publishing it as a
contamination result would itself be the overfitting failure the
discipline layer exists to prevent (logos v0.1.0).

The gate-2 rerun, reconciled to a common consonant-skeleton vocabulary,
yields the honest deliverable: a per-word reproduction table over 40
seeds. \texttt{qwen2.5:72b} reproduces Gordon's \texttt{JA-NE} =
\emph{yn} / ``wine'' in 14/40 seeds and Best 1981's
\texttt{A-SA-SA-RA-ME} = ``oh Asherah!'' in 2/40 --- published readings
the model-free edit-distance baseline never independently reaches (0/40)
--- while never reproducing the obscurer vessel equations
\texttt{KA-RO-PA} and \texttt{SU-PA-RA} (0/40) (logos v0.1.0). Three of
the four tested lexical items --- \texttt{JA-NE}, \texttt{KA-RO-PA},
\texttt{SU-PA-RA} --- are among Gordon's five West-Semitic equations (as
is the separately-treated \texttt{KU-RO} below); the fourth,
\texttt{A-SA-SA-RA-ME}, is Best's, not Gordon's. An external red-team
correction split \texttt{KU-RO}: of its 14/40 hits, 8/40 route through
the Semitic root \emph{kull} and 7/40 are reachable via the
general-knowledge meaning ``total'' (confirmed by tablet arithmetic; the
two categories overlap on one seed), so \texttt{KU-RO} = ``total'' is
\textbf{excluded} as a witness and \texttt{KU-RO} = \emph{kull} is
downweighted (logos v0.1.0). The surviving gradient --- famous readings
reproduced, obscure ones not --- is what C.4 detects:
\textbf{reproduction of published readings weighted by web-prevalence},
not verified correctness. The gate-2 rate itself (mean 0.646) remains
confounded by the baseline's representational ceiling and is not
reported as a contamination measure.

\hypertarget{l_virgin-memorisation-vs.-reasoning}{%
\subsubsection{8.2 L\_virgin: memorisation
vs.~reasoning}\label{l_virgin-memorisation-vs.-reasoning}}

The §C.2 \texttt{L\_virgin} probe shows a sign in real context and asks
whether the model's behaviour on \emph{published} signs generalises to
signs with no published value (which it cannot have memorised) (logos
v0.1.0). Across 40 seeds, \texttt{qwen2.5:72b} assigned a value to the
20 known AB signs in \textasciitilde20/40 seeds but to the 11 virgin
*-series signs in only \textasciitilde1/40 --- it \textbf{abstains
\textasciitilde97\% of the time} on the untouched signs while
confidently valuing the memorisable ones. The apparent known−virgin
consistency gap (0.601, permutation p = 0.0005) is coverage-confounded
and was flagged, not asserted: the matched-\emph{n} robustness check
returned \texttt{no\_power} (0 of 11 virgin signs reached n ≥ 10), so
the rigorous comparison could not run. The interpretable signal is the
abstention asymmetry itself, consistent with regurgitation-dominated
behaviour --- and, per the module's standing rule, \textbf{consistency
is not correctness}: with no ground truth on the virgin signs, this
bounds how the model behaves, never whether it is right (logos v0.1.0).
The full arc cost ≈ \$2.80 in ephemeral GPU rental.

\hypertarget{information-sufficiency-upper-bound-the-csa-learning-curve-measured-8.3-under-revision}{%
\subsubsection{8.3 Information-sufficiency upper bound: the CSA learning
curve (measured) --- §8.3 UNDER
REVISION}\label{information-sufficiency-upper-bound-the-csa-learning-curve-measured-8.3-under-revision}}

\begin{quote}
\textbf{Editorial note (2026-07-01).} This subsection is under active
revision. A correctness bug in the cognate-ID matcher's serial path was
found and fixed (see below); the earlier ``at chance floor'' lower
bounds are \textbf{retracted}, and the corrected accuracy-vs-size curve
is being finalized --- its larger (Linear-A-scale) syllabic points are
still computing. The numbers quoted below are the validated partial
curve as of this draft; the converged full-range curve will replace
them.
\end{quote}

The learning-curve harness reframes the withdrawn unicity-distance
diagnostic --- a toy-model precondition, never a sufficiency or
refutation claim --- as a sharper, answerable question: \emph{is a
Linear-A-scale corpus (539 distinct word-forms; V = 259 sign-types at N
≈ 5,792 parsed syllable-signs) large enough to recover a decipherment IF
a known cognate existed?} We report the design, its reading rules, and
the corrected measured curve here; the larger-size cells for the
syllabic benchmarks --- the analogues most relevant to Linear A --- are
still landing, so the Linear-A-scale verdict is deferred to the
finalized curve.

The protocol downsamples known-answer benchmarks --- Linear B /
Mycenaean Greek as the primary syllabic analog, sourced through DĀMOS,
and further syllabic corpora --- and measures held-out cognate-ID
recovery against a permutation chance floor, using a published
cognate-ID matcher (Tamburini 2025) that reports the Luo et al.~2019
accuracy metric (logos v0.1.0; github.com/papadopouloskyriakos/logos).
Because every benchmark is handed a real cognate signal Linear A lacks,
any recovery curve is a \textbf{strict upper bound}, readable only two
ways: at-floor at Linear-A scale ⇒ definitively
information-insufficient; above-chance ⇒ the binding constraint is
identifiability --- cognate-absence and uncorrected multiplicity over a
mixed sign inventory --- not corpus size. Neither branch is ever a
decipherability claim, and neither asserts ``there is enough text.''

\textbf{A retracted result and its correction.} In preparing the full
sweep we found and fixed a correctness bug in the matcher, and we
retract the reduced-budget numbers reported in an earlier draft.
Tamburini's \emph{serial} annealing update (\texttt{processes=1}, the
\texttt{\_\_update\_state\_no\_par} path) mutates one shared annealer
state in place with no per-worker isolation, collapsing held-out
accuracy to near chance; only the \emph{parallel} path
(\texttt{processes\ \textgreater{}\ 1}, forked worker processes with
isolated state) is correct. A known-answer check pins this down:
Luvian-Hittite at its full size (59 pairs) reproduces Tamburini's
published Table-3 accuracy of 0.475 \textbf{only} on the parallel path
(0.462 at 2,000 steps, mean of four seeds; per-seed energy ≈ 37,
matching the published reproduction), while the serial path returns
0.0--0.017 (energy ≈ 48--54). Our earlier batched-CUDA energy had been
validated ``bit-identical to the serial reference''; because that
reference was the broken path, \textbf{we retract that validation and
the six reduced-budget ``lower bounds''} (formerly
\texttt{results/csa/curve\_lowerbounds.json}), which reported every
point ``at or within noise of the analytic chance floor.'' That at-floor
characterization was a serial-path artifact --- not under-convergence,
and not information-insufficiency.

\textbf{The corrected curve splits informatively by benchmark type.}
Re-run on the parallel path (\texttt{processes\ =\ 4}, CPU, 2,000 steps,
four seeds each; analytic chance floors 0.013--0.05), the
\emph{non-syllabic} benchmarks recover far above floor even at small
downsample sizes --- phoenician-Ugaritic 0.66--0.74 across sizes 21--79,
Luvian-Hittite rising 0.28 → 0.43 → 0.46 across sizes 30--59 (its
full-size 0.46 matching the published 0.475). The \emph{syllabic}
benchmarks --- Linear-B/Greek and Cypriot/Greek, the closest analogues
to syllabic Linear A --- by contrast remain \textbf{at floor at these
small downsample sizes} (Linear-B/Greek ≈ 0.01 at size 92; Cypriot/Greek
≈ 0.08 at size 69). On the syllabic case cognate-ID recovery evidently
needs a larger corpus, and those larger cells (up toward Linear-A scale,
≈ 540 word-forms) are still computing. \textbf{The Linear-A-relevant
sufficiency verdict therefore awaits the larger syllabic sizes, and this
subsection is marked under revision pending them.} What is already
established and does not depend on the remaining cells: the matcher is
\emph{correct} (it reproduces the published known-answer accuracy on the
parallel path), the earlier ``uninformative, at-floor'' state is
\emph{withdrawn}, and the two-branch reading of the curve above is the
frame in which the finalized result will be reported. Per-step cost
grows roughly quadratically in lexicon size, which is why the syllabic
Linear-A-scale points are the slow, still-completing tail. The
phonological / sound path, by contrast, is no longer data-deferred: the
DĀMOS Mycenaean holdings are now harvested and ingested, and §7.3
reports that even the full corpus leaves the cross-script distributional
alignment at chance --- so that bound is a property of the signal, not a
promissory note against data yet to arrive. We publish the corrected
harness, per-size chance floor, and the (regenerating) curve so it is
\textbf{reproduced and audited rather than asserted} (concept DOI
10.5281/zenodo.21087572).

\hypertarget{discussion-limitations-and-conclusion}{%
\subsection{9. Discussion, limitations, and
conclusion}\label{discussion-limitations-and-conclusion}}

\hypertarget{the-discipline-is-the-contribution}{%
\subsubsection{9.1 The discipline is the
contribution}\label{the-discipline-is-the-contribution}}

The deliverable of this work is not a reading of Linear A. It is a
\emph{harness} --- a literature index with a known-reading versus
literature-unseen partition, a fabricated-language control corpus,
model-ablation contamination checks, multiple-comparison correction, and
pre-registration, all graded mechanically from persisted artifacts ---
together with a body of pre-registered, calibrated nulls, each with its
borderline case named. The stratified morphology re-run is the template:
at the pre-registered α = 0.01, no pre-registered affix is cross-stratum
stable above the bigram floor, robust to the genre-bucketing choice; the
closest candidates are the Duhoux/Valério H3 nominal affixes
(\texttt{i-} ``to/at'', \texttt{-te/-ti} ``from/of''), which surface
only in the doubly-relaxed corner (α = 0.05 \emph{and} loose bucketing)
(logos v0.1.0). A null that names its threshold, its sensitivity, and
its near-miss is a stronger referee object than a flat negative.

The axis this harness targets is the \textbf{false-positive /
multiplicity axis} --- the mechanism by which a determined analyst
``translates'' almost anything on a small corpus. As set out in the
Introduction, this is a property of inference methodology and
epistemics, not of the object or its data; it is complementary to, not a
member of, the fifteen data-and-object obstacles enumerated by Braović
et al.~(2024), and distinct from the evaluation problem their §6.3
gestures at.

\hypertarget{two-named-failure-modes-for-the-methods-literature}{%
\subsubsection{9.2 Two named failure modes for the methods
literature}\label{two-named-failure-modes-for-the-methods-literature}}

\textbf{(i) Narrative capture.} Structure says X; the label says the
desired Y; the label wins. In the June-2026 *301 claim the published
table segments an agglutinative, prefix-stacking verb, yet the assigned
language family is Semitic and the desired output (``prayers to a
Goddess in a matriarchal society'') operates as a fixed prior, so the
conclusion conforms to the prior rather than emerging from constrained
readings. This is distinct from statistical overfitting and must be
named separately. Notably, the current synthesis independently concludes
Minoan is agglutinative and a language isolate (Salgarella 2025 §8),
which the segmentation matches while the Semitic label does not --- the
mismatch is empirical, between the label and the primary descriptions of
the form, not an a-priori typological impossibility.

\textbf{(ii) Free parameters exceeding the information floor.} A lexicon
reported as ``408 terms and counting'' on a fixed corpus has more
degrees of freedom than the data can identify; a decipherment whose
lexicon grows weekly is underdetermined by construction. The
minimum-description-length / unicity \emph{diagnostic}
(\texttt{k\ ≤\ U\_floor}) is built to flag exactly this --- after
pre-submission review it is \textbf{reported next to every claim rather
than enforced as a graduation clause} (§3), because clearing it does not
certify identifiability: the single sign the whole reading pivots on
carries no cross-script anchor (its AB-number and Linear B columns are
empty), so its value was assigned to fit the desired root --- the
textbook free-parameter case the machinery flags as underdetermined. The
refutation rests on the claim's own published table, archived in the
logos artifact (v0.1.0).

\hypertarget{the-inverted-pipeline}{%
\subsubsection{9.3 The inverted pipeline}\label{the-inverted-pipeline}}

These failure modes trace to sequencing. The *301 claim publicly commits
to an inverted pipeline: declare ``officially deciphered,'' then vet
later. logos's documented inverse is \textbf{pre-register → held-out →
locked artifact}: a hypothesis is hash-keyed and timestamped before it
is tested, graded mechanically against held-out material, and the
artifact is fixed before any claim. The libation-formula corpus
(\textasciitilde50 inscribed objects over \textgreater5 findspots,
spanning peak sanctuaries and caves) offers a natural held-out
cross-validation, but a leave-one-site-out split over so few formula
sites has low statistical power --- and site-level CV is in any case
necessary, not sufficient, since within-site scribal-hand correlation
must also be controlled (Salgarella \& Judson 2024). This formula corpus
should not be conflated with the 99 stone-vessel inscriptions of the
``libation'' genre stratum (§6): they are different objects.

\hypertarget{identifiability-not-corpus-size-is-the-binding-constraint}{%
\subsubsection{9.4 Identifiability, not corpus size, is the binding
constraint}\label{identifiability-not-corpus-size-is-the-binding-constraint}}

We take care not to write a defeatist framing, and equally not to
over-claim from corpus size. On the real corpus (1,341 inscriptions; N ≈
5,792 parsed syllable-signs, V = 259) the toy-model unicity distance is
U ≈ 204--415 signs --- a \emph{precondition} diagnostic only. We
explicitly withdraw any reading of this figure as a sufficiency or
refutation claim: a lossy, non-injective sign→sound channel may have no
unicity distance at all, so ``\(U \ll N\)'' does \textbf{not} license
the statement that there is enough text to pin a decipherment (this is
the stance of §4, and §9 now aligns with it). The binding constraint is
\emph{identifiability}: the absence of a known cognate language to map
to, an uncorrected multiple-testing search burden, and a mixed sign
inventory. (Note that the \textasciitilde7,400 sign-occurrence figure
sometimes quoted --- Schoep 2002; Salgarella 2025 §1--3 --- counts a
different channel than the 5,792 parsed syllable-signs used for U here.)
Salgarella (2025 §8) reaches the field-level version of this by a
different route: the etymological method is ``fraught with problems,''
vocabulary comparison is ``not probative,'' and she names
digital/statistical analysis as ``the next desideratum in the field.''

\hypertarget{study-wide-error-budget}{%
\subsubsection{9.5 Study-wide error
budget}\label{study-wide-error-budget}}

Multiplicity discipline operates at two levels, and we are explicit
about both. Within a single probe, matches are deflated for the signs ×
roots × hypotheses searched (the effective-n machinery of §3--4). Across
the study, the unit of multiplicity is the \textbf{per-probe
pre-registration}: the \textbf{five executed probes} (morphology,
metrology, phonology, cross-script image alignment, and contamination)
each carry their own pre-committed α, and every probe is reported
regardless of outcome. (The sufficiency sweep of §8.3 is a known-answer
\emph{upper-bound} diagnostic on benchmarks handed a real cognate signal
--- not a Linear A probe --- so it is not among these five; on the
corrected matcher it now returns a measured accuracy-vs-size curve, §8.3
under revision, and its above-chance recovery on known-answer benchmarks
is true by construction, never a Linear A positive requiring
multiplicity protection.) There is no selective reporting: the one probe
that returns an above-chance number --- the cross-script image leg
(§7.3) --- is reported as a circular \emph{engineering demonstration},
not a positive. On the battery-level correction the paper's own rule
demands: four of the five executed probes are nulls sitting far from any
threshold (metrology p = 1.0; phonology p = 0.40 / 0.86; morphology at
no power; contamination an artifact by construction), so a
family-wise/FDR correction over the probe-level tests is arithmetically
trivial and only makes each \emph{more} null --- and the fifth probe
surfaces \textbf{no positive for such a correction to protect} (the
image leg is capped at ≤ 0.75 and relabeled circular). Battery-level
multiplicity is therefore not load-bearing here: there is no surviving
positive that a correction could threaten. Nothing in the battery is
read as a decipherment.

\hypertarget{limitations}{%
\subsubsection{9.6 Limitations}\label{limitations}}

The nulls are honestly bounded, not decisive. Full effective-n deflation
of the morphology z-test is deferred; the current control is a
conservative ≥ 2-distinct-sites presence gate, which would only
strengthen the null (logos v0.1.0). No held-out \emph{decipherment} has
been achieved, and the study reports \textbf{no offensive positive} ---
the deliverable is calibrated absence, not a reading. The one probe that
recovers an above-chance number, cross-script sign-image alignment
(direct-NN 0.410 {[}0.18, 0.64{]}; Procrustes 0.185 {[}0.00, 0.36{]}),
is a font-robust but circular \emph{engineering demonstration} (§7.3): a
font-swap control refutes the typeface-artifact worry, but the A↔B
anchor values were themselves assigned by sign shape, so the recovery
re-derives a shape-defined mapping and is truth-capped at ≤ 0.75, never
a reading. No phonetic claim is made, and no sign claim should be
advanced without review by Aegean epigraphers. Chronological
cross-validation is largely foreclosed because roughly 74\% of the
corpus sits in the single LM IB horizon, leaving site-held-out splits as
the only realistic partition. And the corpus itself --- small, partly
lossy, with no bilingual yet found (Salgarella 2025 §8) --- sets the
ceiling.

\hypertarget{conclusion-and-reproducibility}{%
\subsubsection{9.7 Conclusion and
reproducibility}\label{conclusion-and-reproducibility}}

If a Linear A decipherment is attainable on present evidence, this
framework is built to recognize it; the aim is a real reading, and the
calibrated null is the field's insurance policy against another false
one. The code, pre-registrations, and all findings are open under the
MIT licence and archived on Zenodo (concept DOI 10.5281/zenodo.21087572;
repository https://github.com/papadopouloskyriakos/logos). The licensed
epigraphic corpora --- GORILA-, SigLA-, lineara.xyz-, and DĀMOS-derived
--- are not redistributed, per their CC BY-NC-SA / all-rights-reserved
terms; each source's licence and retrieval path is documented
(data-provenance), so the harness is reproducible against a locally
obtained corpus. This work is independent and not affiliated with, or
endorsed by, any of the cited projects or their authors. \#\# References

\emph{Bibliographic details below are the works cited above; the full,
continuously-maintained bibliography (including works consulted) lives
in the repository at \texttt{docs/references.md}. Per-source digests are
in \texttt{docs/related/}. Works cited only through a secondary source
are marked ``secondary citation'' at the point of use.}

\begin{itemize}
\tightlist
\item
  Bailey, D. H., \& López de Prado, M. (2014). The Deflated Sharpe
  Ratio: Correcting for Selection Bias, Backtest Overfitting, and
  Non-Normality. \emph{Journal of Portfolio Management} 40(5):94--107.
\item
  Bennett, E. L., Jr.~(1950). Fractional Quantities in Minoan
  Bookkeeping. \emph{American Journal of Archaeology} 54(3):204--222.
\item
  Berg-Kirkpatrick, T., \& Klein, D. (2011). Simple Effective
  Decipherment via Combinatorial Optimization. \emph{Proc. EMNLP 2011},
  313--321.
\item
  Best, J. G. P. (1981). YASSARAM! \emph{Talanta} 13:17--21.
\item
  Braović, M., Krstinić, D., Štula, M., \& Ivanda, A. (2024). A
  Systematic Review of Computational Approaches to Deciphering Bronze
  Age Aegean and Cypriot Scripts. \emph{Computational Linguistics}
  50(2):725--779. doi:10.1162/coli\_a\_00514.
\item
  Corazza, M., Ferrara, S., Montecchi, B., Tamburini, F., \& Valério, M.
  (2021). The mathematical values of fraction signs in the Linear A
  script. \emph{Journal of Archaeological Science} 125:105214.
  doi:10.1016/j.jas.2020.105214.
\item
  Corazza, M., et al.~(2022). Unsupervised deep learning supports
  reclassification of the Bronze Age Cypriot writing system
  (\emph{Sign2Vec\_d}). \emph{PLoS ONE} 17(7):e0269544. DOI
  \href{https://doi.org/10.1371/journal.pone.0269544}{10.1371/journal.pone.0269544}.
\item
  Daggumati, S., \& Revesz, P. Z. (2019). Data mining ancient scripts to
  investigate their relationships. \emph{Proceedings of the 23rd
  International Database Applications \& Engineering Symposium
  (IDEAS/DEXA)}. And (2023), \emph{Information} 14(4):227.
\item
  Davis, B. (2013). Syntax in Linear A: The Word-Order of the `Libation
  Formula'. \emph{Kadmos} 52(1):35--52. doi:10.1515/kadmos-2013-0003.
\item
  Davis, B. (2014). \emph{Minoan Stone Vessels with Linear A
  Inscriptions.} Aegaeum 36. Leuven: Peeters.
\item
  Di Mino, D. (2026). ``Linear A cracked'': A-TA-I-*301-WA-JA as
  West-Semitic N-W-Y (the June-2026 announcement analysed in §5). Online
  announcement --- no formal venue; analysed here as a public claim,
  with the timestamped reference copy of the claim and its table carried
  as the archived record in the logos artifact (v0.1.0).
\item
  Duhoux, Y. (1978). On the (agglutinative) typology of the Minoan
  language (secondary citation, as characterised in Salgarella 2025 §8).
\item
  Duhoux, Y. (1992). Variations morphosyntaxiques dans les textes votifs
  linéaires A. \emph{Cretan Studies} 3:65--88. (Source of the i-*301 =
  ``give/dedicate'' structural gloss, adopted by Davis 2013 and endorsed
  by Salgarella 2025 §7.2.)
\item
  Duhoux, Y. (1997). The Linear A locative/allative prefix i-/j-
  ``to/at'' (secondary citation, as cited in Salgarella 2025 §8).
\item
  Ferrara, S., \& Tamburini, F. (2022). Advanced techniques for the
  decipherment of ancient scripts. \emph{Lingue e Linguaggio}
  21(2):239--259. DOI
  \href{https://doi.org/10.1418/105964}{10.1418/105964}.
\item
  Ferrara, S., Montecchi, B., \& Valério, M. (2021). What is the
  `Archanes Formula'? Deconstructing and reconstructing the earliest
  attestation of writing in the Aegean. \emph{Annual of the British
  School at Athens} 116:43--62. DOI
  \href{https://doi.org/10.1017/S0068245420000155}{10.1017/S0068245420000155}.
\item
  Fuls, A. (2015). Classifying undeciphered writing systems.
  \emph{Historische Sprachforschung} 128(1):42--58.
\item
  Gelb, I. J., \& Whiting, R. M. (1975). Methods of decipherment.
  Secondary citation: the Type 0--III script typology is used here only
  as characterised in Ferrara and Tamburini (2022, §3.5), not from the
  original, which we do not cite firsthand.
\item
  Godart, L., \& Olivier, J.-P. (1976--1985). \emph{Recueil des
  inscriptions en Linéaire A} (GORILA). Études crétoises 21. Paris:
  Geuthner.
\item
  Gordon, C. H. (1957). Notes on Minoan Linear A. \emph{Antiquity}
  31(124):237--240.
\item
  Gordon, C. H. (1966). \emph{Evidence for the Minoan Language.}
  Ventnor, NJ: Ventnor Publishers.
\item
  Loh, C. J. S., \& Perono Cacciafoco, F. (2020). A New Approach to the
  Decipherment of Linear A, Stage 2: \ldots{} a `Brute Force Attack' on
  an Undeciphered Writing System. \emph{Grapholinguistics in the 21st
  Century 2020}, 927--943. doi:10.36824/2020-graf-cacc.
\item
  Luo, J., Cao, Y., \& Barzilay, R. (2019). Neural Decipherment via
  Minimum-Cost Flow: From Ugaritic to Linear B. \emph{Proc. ACL 2019},
  3146--3155.
\item
  Luo, J., Hartmann, F., Santus, E., Barzilay, R., \& Cao, Y. (2021).
  Deciphering Undersegmented Ancient Scripts Using Phonetic Prior.
  \emph{Transactions of the ACL} 9:69--81.
\item
  Eu, N., Xu, D., \& Perono Cacciafoco, F. (2019). Coding to decipher
  Linear A: consonantal comparison against time-compatible lexical
  lists. \emph{Proceedings of the 2019 Pacific Neighbourhood Consortium
  (PNC) Annual Conference and Joint Meetings}, Singapore.
\item
  Packard, D. W. (1974). \emph{Minoan Linear A.} Berkeley/Los Angeles:
  University of California Press. (Constructs nine deliberately
  fabricated control decipherments; the genuine Ventris-value transfer
  beats their average by \textasciitilde2:1, improving to
  \textasciitilde3:1 when context-adjusted.)
\item
  Rendsburg, G. A. (1996). ``Someone Will Succeed in Deciphering
  Minoan'': Cyrus H. Gordon and Minoan Linear A. \emph{Biblical
  Archaeologist} 59(1):36--43.
\item
  Salgarella, E. (2019). Drawing lines: The palaeography of Linear A and
  Linear B. \emph{Kadmos} 58(1/2):61--92.
\item
  Salgarella, E. (2020). \emph{Aegean Linear Script(s): Redefining the
  Relationship between Linear A and Linear B.} Cambridge: Cambridge
  University Press.
\item
  Salgarella, E. (2025). \emph{Writing in Bronze Age Crete: ``Minoan''
  Linear A.} Cambridge Elements in Writing in the Ancient World.
  Cambridge University Press. doi:10.1017/9781009520041.
\item
  Salgarella, E., \& Judson, A. P. (2024). Signs of the times? Testing
  the chronological significance of Linear A and B palaeography. In
  \emph{KO-RO-NO-WE-SA} (Ariadne Suppl. 5), 359--379.
\item
  Schoep, I. (2002). \emph{The administration of neopalatial Crete: a
  critical assessment of the Linear A tablets.} Minos Suppl. 17.
\item
  Schrijver, P. (2014). Fractions and food rations in Linear A.
  \emph{Kadmos} 53(1-2):1--44. doi:10.1515/kadmos-2014-0001.
\item
  Snyder, B., Barzilay, R., \& Knight, K. (2010). A Statistical Model
  for Lost Language Decipherment. \emph{Proc. ACL 2010}, 1048--1057.
\item
  Sommerschield, T., Assael, Y., Pavlopoulos, J., et al.~(2023). Machine
  Learning for Ancient Languages: A Survey. \emph{Computational
  Linguistics} 49(3):703--747.
\item
  Steele, P. M., \& Meissner, T. (2017). From Linear B to Linear A: the
  problem of the backward projection of sound values. In P. M. Steele
  (ed.), \emph{Understanding Relations Between Scripts: The Aegean
  Writing Systems}, 93--110. Oxford: Oxbow.
\item
  Tamburini, F. (2025). On automatic decipherment of lost ancient
  scripts relying on combinatorial optimisation and coupled simulated
  annealing (CSA\_OptMatcher), evaluated by the Luo et al.~(2019)
  accuracy metric. \emph{Frontiers in Artificial Intelligence}
  8:1581129. DOI
  \href{https://doi.org/10.3389/frai.2025.1581129}{10.3389/frai.2025.1581129}.
  Code: github.com/ftamburin/CSA\_OptMatcher.
\item
  Thomas, R. (2020). Some reflections on morphology in the language of
  the Linear A libation formula. \emph{Kadmos} 59(1--2):1--23. DOI
  \href{https://doi.org/10.1515/kadmos-2020-0001}{10.1515/kadmos-2020-0001}.
\item
  Valério, M. (2007). Analysis of Linear A suffixation (secondary
  citation, as cited in Salgarella 2025 §8).
\item
  Ventris, M., \& Chadwick, J. (1953). Evidence for Greek Dialect in the
  Mycenaean Archives. \emph{Journal of Hellenic Studies} 73:84--103.
\end{itemize}

\begin{center}\rule{0.5\linewidth}{0.5pt}\end{center}

\emph{Software \& data.} logos (v0.1.0). Zenodo. Concept DOI
\href{https://doi.org/10.5281/zenodo.21087572}{10.5281/zenodo.21087572};
repository \url{https://github.com/papadopouloskyriakos/logos}. Licensed
corpora (GORILA-, SigLA-, lineara.xyz-, DĀMOS-derived) are not
redistributed; see \texttt{docs/data-provenance.md}.

\end{document}
